\begin{lemma}
\label{editorial-source-lemma-field-finite-type-over-domain}
Let $R$ be a ring. Let $K$ be a field.
If $R \subset K$ and $K$ is of finite type over $R$,
then there exists an $f \in R$ such that $R_f$ is a field,
and $K/R_f$ is a finite field extension.
\end{lemma}

\begin{proof}
By Lemma \ref{editorial-source-lemma-characterize-image-finite-type} there
exist a nonempty open $U \subset \Spec(R)$
contained in the image $\{(0)\}$ of $\Spec(K) \to \Spec(R)$.
Choose $f \in R$, $f \not = 0$ such that $D(f) \subset U$, i.e.,
$D(f)=\{(0)\}$. The localization $R_f$ is a domain whose only
prime ideal is $(0)$. Every nonunit belongs to a maximal ideal,
so its only nonunit is zero and $R_f$ is a field. The original
inclusion $R\subset K$ extends to $R_f\subset K$, and the original
finite set of $R$-algebra generators of $K$ is also a set of
$R_f$-algebra generators. The Hilbert Nullstellensatz
(Theorem \ref{editorial-source-theorem-nullstellensatz}) therefore makes $K/R_f$ finite.
In fact the canonical inclusion $R_f\subset\operatorname{Frac}(R)$
is an equality: for $r/s\in\operatorname{Frac}(R)$ with $s\ne0$,
the element $s/1$ is nonzero in the field $R_f$ and has an inverse
there, so $r/s$ belongs to $R_f$. Thus the same $f$ realizes the
original fraction field as a principal localization of $R$.
\end{proof}